\documentclass[../Main.tex]{subfiles}

\begin{document}

\label{chap:code}

\begin{mmaCell}[morepattern={x_, x}]{Input}
SetDirectory@NotebookDirectory[];
\end{mmaCell}

	\mmaSec{Importación de datos}

\begin{mmaCell}[morepattern={T_,P_}]{Input}
experimentalData=Import["experimental.txt","Table"]/.
   \{\mmaPat{θ_},P_\}->\{\mmaUnd{θ},P/\mmaSup{10}{6}\};
\end{mmaCell}

	\mmaSec{Definición de constantes y parámetros globales}
 
\begin{mmaCell}[moredefined={nAir, rDB}]{Input}
\mmaDef{λRb}=\mmaUnit{780.241nm};
nAir=1.000277;
rDB=0.85(*comentario en Mathematica*);
\end{mmaCell}

\begin{mmaCell}[addtoindex=2]{Input}
diode=\{\mmaUnd{μ}->\mmaUnit{782.2nm}, 
   \mmaUnd{σ}->\mmaUnit{10nm},
   L->\mmaUnit{0.282mm},
   n->3.5,
   r1->rDB, r2->0.15\};
cavity=\{L->\mmaUnit{50mm},
   n->nAir, r1->rDB, r2->0.3\};	
\end{mmaCell}

Texto documentando el código

\begin{mmaCell}[addtoindex=1,moredefined={f}]{Input}
\mmaDef{∆νD}=UnitConvert\Big[\mmaFrac{\mmaUnit{c}}{2 n L},"GHz"\Big]/.diode
\mmaDef{∆νC}=UnitConvert\Big[\mmaFrac{\mmaUnit{c}}{2 n L},"GHz"\Big]/.cavity
\end{mmaCell}

\begin{mmaCell}{Output}
\mmaUnit{151.871 GHz}
\end{mmaCell}

\begin{mmaCell}{Output}
\mmaUnit{2.99709 GHz}
\end{mmaCell}

\end{document}